\pdfbookmark{Общая характеристика работы}{characteristic}             % Закладка pdf
\section*{Общая характеристика работы}

\newcommand{\actuality}{\pdfbookmark[1]{Актуальность}{actuality}\underline{\textbf{\actualityTXT}}}
\newcommand{\progress}{\pdfbookmark[1]{Разработанность темы}{progress}\underline{\textbf{\progressTXT}}}
\newcommand{\aim}{\pdfbookmark[1]{Цели}{aim}\underline{{\textbf\aimTXT}}}
\newcommand{\tasks}{\pdfbookmark[1]{Задачи}{tasks}\underline{\textbf{\tasksTXT}}}
\newcommand{\aimtasks}{\pdfbookmark[1]{Цели и задачи}{aimtasks}\aimtasksTXT}
\newcommand{\novelty}{\pdfbookmark[1]{Научная новизна}{novelty}\underline{\textbf{\noveltyTXT}}}
\newcommand{\influence}{\pdfbookmark[1]{Практическая значимость}{influence}\underline{\textbf{\influenceTXT}}}
\newcommand{\methods}{\pdfbookmark[1]{Методология и методы исследования}{methods}\underline{\textbf{\methodsTXT}}}
\newcommand{\defpositions}{\pdfbookmark[1]{Положения, выносимые на защиту}{defpositions}\underline{\textbf{\defpositionsTXT}}}
\newcommand{\reliability}{\pdfbookmark[1]{Достоверность}{reliability}\underline{\textbf{\reliabilityTXT}}}
\newcommand{\probation}{\pdfbookmark[1]{Апробация}{probation}\underline{\textbf{\probationTXT}}}
\newcommand{\contribution}{\pdfbookmark[1]{Личный вклад}{contribution}\underline{\textbf{\contributionTXT}}}
\newcommand{\publications}{\pdfbookmark[1]{Публикации}{publications}\underline{\textbf{\publicationsTXT}}}

{\actuality} Последние десятилетия характеризуются активным внедрением и распространением беспилотных транспортных средств (БТС) в различных сферах человеческой деятельности. Примеры таких сфер -- логистика, промышленность, сельское хозяйство, горное дело, городская инфраструктура. Эксплуатация БТС четвёртого уровня автоматизации по классификации SAE~J3016 вышла за пределы испытательных полигонов, и с расширением области эксплуатации ужесточаются требования к безопасности их движения. Развитие беспилотного транспорта играет важную роль и в достижении национальных целей и задач развития Российской Федерации. Так, <<Транспортные технологии для различных сфер применения (море, земля, воздух), в том числе беспилотные и автономные системы>> входят в перечень критических технологий Российской Федерации. Транспортная стратегия РФ на период до 2030 года включает организацию движения БТС на автомобильной дороге общего пользования. В рамках Стратегии развития автомобильной промышленности Российской Федерации до 2035 года создана инициатива социально-экономического развития <<Беспилотные логистические коридоры>>, предусматривающая проведение мероприятий по внедрению на автомобильных дорогах Государственной компании интеллектуальных транспортных систем, ориентированных в том числе на обеспечение движения высокоавтоматизированных транспортных средств.

Основой безопасного движения БТС служит подсистема восприятия, строящая по данным бортовых датчиков модель окружающего пространства. К такой модели предъявляются два встречных требования. С одной стороны, она должна быть достаточно точной, чтобы обеспечить решение задач использующими её алгоритмами. С другой -- лишённой избыточной детализации: чем сложнее модель, тем больше времени требует работа с ней, а минимизация задержек в системах автопилотирования критически важна, поскольку любая задержка означает, что решение принимается по устаревшему описанию окружения. Поэтому состав сведений, хранимых в модели, определяется тем, какие задачи по ней решаются. В частности, алгоритмам планирования маршрута требуется модель, описывающая расположение зон, движение в которых возможно, невозможно либо о проходимости которых сведения отсутствуют.

Точность описания каждого из этих классов зон одинаково существенна, а ошибки модели опасны в обе стороны: занижение занимаемой препятствием области ведёт к столкновению, завышение -- к ложным остановкам и снижению транспортной работы, а области, ошибочно отнесённые к свободным при фактическом отсутствии наблюдений, лишают алгоритмы планирования возможности действовать консервативно в условиях окклюзии. При этом ни один тип бортовых датчиков не позволяет построить такую модель в одиночку: каждый ограничен собственным принципом работы -- угловым разрешением, дальностью, чувствительностью к погодным условиям либо отсутствием прямого измерения расстояния. Поэтому современные БТС оснащаются наборами разнородных датчиков, а качество модели определяется тем, каким образом их показания комплексируются в единое непротиворечивое описание окружения.

Задача построения модели окружающего пространства, в том числе и модели проходимости окружающего пространства, в робототехнике широко исследована, однако её решение в условиях ограниченного вычислительного бюджета остаётся предметом исследований. Исследование этой задачи нашло отражение в работах М. Павоне, А. Элфеса, Х. Моравека, С. Труна, Б.М. Миллера, Д.А. Юдина и Л.М.Г. Гонсалвеса. Среди исследований задачи комплексирования данных стоит выделить работы Б. Халеги, Д.П. Николаева, Ю.В. Визильтера.

Наибольший объём сведений об окружающей БТС сцене даёт камера: по угловому разрешению она превосходит остальные бортовые датчики и позволяет получить семантическое описание наблюдаемого -- границы проезжей части, классы объектов, дорожную разметку. Однако изображение задаёт лишь направление на объект и не содержит расстояния до него. Поэтому учёт визуальных данных в модели окружающего пространства требует отдельного этапа -- проецирования результатов обработки изображения, детекций объектов и семантической сегментации, из системы координат изображения на плоскость движения БТС. От точности этого проецирования напрямую зависит корректность положения препятствий в модели.

Существующие способы такого проецирования тяготеют к двум крайностям: высокой точности либо низкой вычислительной стоимости. На одном краю -- методы, опирающиеся на известные калибровки камеры и геометрическое предположение о форме поверхности движения. Распространённое допущение о совпадении этой поверхности с плоскостью нарушается как рельефом местности, так и требованиями к конструкции дорожного покрытия. К этому добавляется отклонение фактической внешней калибровки от измеренной, вызванное загрузкой БТС, работой подвески и смещением датчиков при длительной эксплуатации. Вносимая обоими факторами погрешность растёт с удалением от камеры и на тех дальностях, на которых принимается решение об объезде или остановке, составляет единицы метров. На другом краю -- обучаемые преобразования признаков изображения в вид сверху, не использующие явной модели поверхности. Методы этой группы демонстрируют наивысшее качество проецирования, однако требуют производительного  графического ускорителя для запуска нейросетевой модели, принимающей данные от каждого источника данных; размеченных мультимодальных наборов данных для обучения; а также переобучения при любом изменении состава и расположения датчиков, что для мелкосерийных платформ с уникальной сенсорной конфигурацией практически неприемлемо. Предлагаемы в настоящей работе алгоритмы проецирования визуальных данных занимают промежуточное положение между этими крайностями: предположение о совпадении поверхности движения с плоскостью заменяется оценкой её формы в существенно более широком классе поверхностей, но без привлечения обучаемых преобразований в вид сверху.

Второе ограничение имеет тот же источник -- цену одного дополнительного источника данных, -- но проявляется не при проецировании, а при учёте данных в общей модели. Увеличение числа источников данных -- естественный способ сократить слепые зоны БТС и повысить достоверность модели, однако показания каждого из них должны быть сведены в общую модель, а эта операция выполняется над разделяемым ресурсом и потому в существующих реализациях сериализуется. В результате задержка между регистрацией данных датчиком и их отражением в модели растёт с числом источников. Положение усугубляется тем, что для широкоиспользуемых моделей, описывающих пространство в виде двумерного массива -- сеточных представлений, требуемый размер массива определяется тормозным путём и растёт как четвёртая степень рабочей скорости БТС, вследствие чего предельное число одновременно учитываемых источников с ростом скорости сокращается столь же быстро. Известные способы распараллеливания построения сеточных моделей применяются на других уровнях -- внутри одного обновления карты на графическом ускорителе либо между независимыми слоями, сведение которых в общую сетку всё равно выполняется последовательно, -- а неблокирующие структуры данных к разделяемой сетке локальной карты, следующей за БТС и потому требующей операций сдвига и приведения всех элементов к общему моменту времени, не применялись.

Таким образом, качество и актуальность модели, описывающей проходимость окружающего БТС пространства, ограничены двумя факторами: точностью проецирования визуальных данных на плоскость движения без опоры на предположение о плоской поверхности и в пределах вычислительного бюджета бортового вычислителя, а также ростом задержки учёта данных при увеличении числа источников. Оба ограничения имеют общий корень: вычислительная стоимость известных решений растёт с числом учитываемых источников данных, тогда как ресурсы бортового вычислителя разделяются между всеми подсистемами БТС. Разработка вычислительно эффективных алгоритмов комплексирования разнородных данных, а также структуры данных и алгоритма её обновления, допускающих одновременный учёт показаний нескольких источников, представляет собой актуальную научно-техническую задачу.

{\aim} работы является снижение задержки построения и повышение точности модели проходимости окружающего пространства БТС при одновременном учёте данных нескольких разнородных источников в условиях ограниченных вычислительных ресурсов бортового вычислителя.

Для~достижения поставленной цели необходимо было решить следующие {\tasks}:
\begin{enumerate}[beginpenalty=10000] % https://tex.stackexchange.com/a/476052/104425
  \item Провести обзор существующих моделей окружающего пространства БТС и алгоритмов их построения для выбора базовой модели, описывающей проходимость пространства, и алгоритма её построения, которые будут использованы при разработке дальнейших алгоритмов.
  \item Разработать алгоритмы оценки положения объектов в окружающем пространстве БТС, а также областей проезжей
  части, доступных для движения БТС, на основе визуальных и лидарных данных, предназначенные для работы на бортовом вычислителе БТС в режиме реального времени.
  \item Разработать программный комплекс, обеспечивающий построение модели окружающего пространства БТС на основе комплексирования разнородных данных с бортовых датчиков БТС.
  \item Разработать структуру данных для представления модели окружающего пространства и параллельный алгоритм её обновления несколькими источниками данных.
\end{enumerate}


{\novelty}
\begin{enumerate}[beginpenalty=10000] % https://tex.stackexchange.com/a/476052/104425
  \item Предложен алгоритм оценки положения окружающих ТС на плоскости движения БТС по данным монокулярной камеры, отличающийся от известных алгоритмов вписывания в ориентированную ограничивающую рамку тем, что в его основу положена модель L-образной формы видимой границы ТС, ранее применявшаяся лишь к кластерам точек дальномерных датчиков.
  \item Предложен алгоритм оценки положения окружающих ТС на плоскости движения БТС по комплексированию визуальных и лидарных данных, отличающийся от известных тем, что сегментация лидарного облака на точки земли и точки препятствий выполняется без предположения о плоской либо кусочно-плоской поверхности движения -- на основе оценки высоты поверхности регрессией на гауссовских процессах, -- что повышает точность оценки положения во всём диапазоне рассмотренных дальностей.
  \item Предложен алгоритм проекции семантической сегментации проезжей части с изображения на плоскость движения БТС, отличающийся от известных тем, что проецирование выполняется не на фиксированную плоскость, а на оценённую по лидарным данным поверхность движения в проецируемой области, что снижает ошибку описания свободного пространства на строящейся карте проходимости.
  \item Впервые количественно показано, что для класса алгоритмов построения карты проходимости пространства, не использующих обучаемых преобразований в вид сверху и применимых на бортовом вычислителе БТС, учёт визуальных данных с камер повышает качество карты по сравнению с картой, строящейся исключительно по лидарным данным, причём улучшение проявляется в том числе по метрике, учитывающей задачу планирования пути (PFC-MSE).
  \item Предложены структура данных для представления карты проходимости пространства и алгоритм её обновления, отличающиеся от известных реализаций тем, что операции над разделяемой сеткой выполняются без использования примитивов взаимного исключения, что позволяет обновлять карту нескольким источникам данных одновременно и снижает время задержки между моментом регистрации данных датчиком и моментом обновления карты.
\end{enumerate}

{\influence}. Разработанные алгоритмы и комплекс технических средств, обеспечивающие построение модели окружающего пространства БТС,  реализованные в данной работе, внедрены в беспилотные транспортные средства <<EVO.1>> и <<EVO.3>>, разработанные компанией ООО <<ЭвоКарго>>. Внедрение разработанных в данной диссертации методов подтверждено соот­ветствующими актами. На результаты, полученные в работе, получены свидетельства о государственной регистрации программ для ЭВМ:

\begin{enumerate}
	\item №~2025617212 Программное обеспечение <<Система восприятия N1 V3>>.
	\item №~2025669744 Программное обеспечение <<Способ построения карты проходимости в окрестности высокоавтоматизированного беспилотного грузового транспортного средства>>.
\end{enumerate}

На основные результаты работы получен патент №~2853062 <<Способ построения карты проходимости в окрестности высокоавтоматизированного беспилотного грузового транспортного средства>>.

{\methods} В диссертации используются методы вероятностной робототехники, методы байесовской оценки, вычислительной оптимизации, компьютерного зрения, анализ вычислительной сложности, а также численный эксперимент.

{\defpositions}
\begin{enumerate}[beginpenalty=10000] % https://tex.stackexchange.com/a/476052/104425
  \item На записях показаний датчиков реальных БТС предложенный алгоритм вписывания проекции объекта на виде сверху в ориентированную ограничивающую рамку, опирающийся на модель L-образной формы видимой границы ТС, обеспечивает прирост среднего коэффициента Жаккара оценки положения окружающих ТС относительно сравниваемых общеизвестных алгоритмов вписывания, не использующих обучаемые методы и применённых к тем же входным данным: $0,337$ против $0,328$ у лучшего из них (прирост $2,7$~\%).
  \item Среди алгоритмов оценки положения объектов на плоскости движения БТС, применимых на бортовом вычислителе БТС в режиме реального времени, предложенный алгоритм комплексирования визуальных и лидарных данных, не использующий приближение поверхности дороги плоскостью, обеспечивает наибольший средний коэффициент Жаккара для ТС находящихся на расстоянии до 60 м: прирост $10,0$~\% по сравнению с алгоритмом, использующим только лидарные данные ($0,439$ против $0,399$), и прирост $30,3$~\% по сравнению с алгоритмом, использующим только визуальные данные ($0,439$ против $0,337$).
  \item При построении карты проходимости пространства размером $600\times600$ элементов с разрешением $0,1$~м по данным трёх бортовых камер и двух лидаров с ограничением дальности $40$~м предложенный алгоритм проекции семантической сегментации проезжей части на плоскость движения БТС на основе комплексирования визуальных и лидарных данных в сочетании с базовым лидарным алгоритмом повышает качество карты по сравнению с картой, строящейся исключительно по лидарным данным: среднеквадратичная ошибка относительно эталонной карты снижается в среднем с $0,231$ до $0,218$ (на $5,6$~\%), а среднеквадратичная ошибка стоимости поиска пути (PathFinding Cost Mean Squared Error, PFC-MSE) -- с $1019,3$ до $924,4$ (на $9,3$~\%), причём по метрике PFC-MSE улучшение достигается во всех пяти рассмотренных сценариях движения.
  \item Предложенные структура данных для представления карты проходимости пространства и алгоритм, допускающий неблокирующее выполнение операций над картой, уменьшают коэффициент линейной аппроксимации зависимости времени одной итерации цикла обновления карты $8$-ю источника данных от размера карты в $13,9$ раза по сравнению с базовой реализацией, использующей блокировки при выполнении операций.
  \item Разработанный программный комплекс построения карты проходимости пространства позволяет подключать новый источник данных, в том числе датчик нового типа, не изменяя код построения карты. Число одновременно подключённых источников масштабируемо и ограничено только вычислительными возможностями бортового вычислителя.
\end{enumerate}

{\reliability} полученных результатов обеспечивается использованием строгого математического аппарата и методов математической статистики. Помимо этого, достоверность подтверждается результатами вычислительных и натурных экспериментов и внедрением разработанных методов в БТС, находящиеся в эксплуатации.

Экспериментальная оценка выполнена на собственных наборах данных, зарегистрированных с реальных БТС. Публичные наборы данных, применяемые для оценки алгоритмов трёхмерной детекции и сегментации (KITTI, SemanticKITTI, nuScenes), для решаемых в работе задач неприменимы: они не содержат ни конфигурации датчиков, соответствующей рассматриваемым БТС, ни разметки проходимости пространства, необходимой для формирования эталонных карт.

{\probation}
Основные результаты работы докладывались~на международных конференциях:
\begin{itemize}
	\item 16th International Conference on Machine Vision (ICMV~2023, Ереван, Армения);
	\item 40th ECMS International Conference on Modelling and Simulation (ECMS~2026, Гримстад, Норвегия).
\end{itemize}
 

{\contribution} Все результаты диссертации, вынесенные на защиту, получены автором самостоятельно. Постановка задач и обсуждение результатов проводились совместно с научным руководителем.

\ifnumequal{\value{bibliosel}}{0}
{%%% Встроенная реализация с загрузкой файла через движок bibtex8. (При желании, внутри можно использовать обычные ссылки, наподобие `\cite{vakbib1,vakbib2}`).
    {\publications} Основные результаты по теме диссертации изложены
    в~XX~печатных изданиях,
    X из которых изданы в журналах, рекомендованных ВАК,
    X "--- в тезисах докладов.
}%
{%%% Реализация пакетом biblatex через движок biber
    \begin{refsection}[bl-author, bl-registered]
        % Это refsection=1.
        % Процитированные здесь работы:
        %  * подсчитываются, для автоматического составления фразы "Основные результаты ..."
        %  * попадают в авторскую библиографию, при usefootcite==0 и стиле `\insertbiblioauthor` или `\insertbiblioauthorgrouped`
        %  * нумеруются там в зависимости от порядка команд `\printbibliography` в этом разделе.
        %  * при использовании `\insertbiblioauthorgrouped`, порядок команд `\printbibliography` в нём должен быть тем же (см. biblio/biblatex.tex)
        %
        % Невидимый библиографический список для подсчёта количества публикаций:
        \phantom{\printbibliography[heading=nobibheading, section=1, env=countauthorvak,          keyword=biblioauthorvak]%
        \printbibliography[heading=nobibheading, section=1, env=countauthorwos,          keyword=biblioauthorwos]%
        \printbibliography[heading=nobibheading, section=1, env=countauthorscopus,       keyword=biblioauthorscopus]%
        \printbibliography[heading=nobibheading, section=1, env=countauthorconf,         keyword=biblioauthorconf]%
        \printbibliography[heading=nobibheading, section=1, env=countauthorother,        keyword=biblioauthorother]%
        \printbibliography[heading=nobibheading, section=1, env=countregistered,         keyword=biblioregistered]%
        \printbibliography[heading=nobibheading, section=1, env=countauthorpatent,       keyword=biblioauthorpatent]%
        \printbibliography[heading=nobibheading, section=1, env=countauthorprogram,      keyword=biblioauthorprogram]%
        \printbibliography[heading=nobibheading, section=1, env=countauthor,             keyword=biblioauthor]%
        \printbibliography[heading=nobibheading, section=1, env=countauthorvakscopuswos, filter=vakscopuswos]%
        \printbibliography[heading=nobibheading, section=1, env=countauthorscopuswos,    filter=scopuswos]}%
        %
        \nocite{*}%
        %
        {\publications} Основные результаты по теме диссертации изложены в~\arabic{citeauthor}~печатных изданиях,
        \ifnum \value{citeauthorvak}>0\relax
        \arabic{citeauthorvak} из которых изданы в журналах, рекомендованных ВАК%
        \fi%
        \ifnum \value{citeauthorscopuswos}>0%
            , \arabic{citeauthorscopuswos} "--- в~изданиях, индексируемых Web of~Science и Scopus%
        \fi%
        \ifnum \value{citeauthorconf}>0%
            , \arabic{citeauthorconf} "--- в~тезисах докладов.
        \else%
            .
        \fi%
        \ifnum \value{citeregistered}=1%
            \ifnum \value{citeauthorpatent}=1%
                Зарегистрирован \arabic{citeauthorpatent} патент.
            \fi%
            \ifnum \value{citeauthorprogram}=1%
                Зарегистрирована \arabic{citeauthorprogram} программа для ЭВМ.
            \fi%
        \fi%
        \ifnum \value{citeregistered}>1%
            Зарегистрированы\ %
            \ifnum \value{citeauthorpatent}>0%
            \formbytotal{citeauthorpatent}{патент}{}{а}{}%
            \ifnum \value{citeauthorprogram}=0 . \else \ и~\fi%
            \fi%
            \ifnum \value{citeauthorprogram}>0%
            \formbytotal{citeauthorprogram}{программ}{а}{ы}{} для ЭВМ.
            \fi%
        \fi%
        % К публикациям, в которых излагаются основные научные результаты диссертации на соискание учёной
        % степени, в рецензируемых изданиях приравниваются патенты на изобретения, патенты (свидетельства) на
        % полезную модель, патенты на промышленный образец, патенты на селекционные достижения, свидетельства
        % на программу для электронных вычислительных машин, базу данных, топологию интегральных микросхем,
        % зарегистрированные в установленном порядке.(в ред. Постановления Правительства РФ от 21.04.2016 N 335)
    \end{refsection}%
    \begin{refsection}[bl-author, bl-registered]
        % Это refsection=2.
        % Процитированные здесь работы:
        %  * попадают в авторскую библиографию, при usefootcite==0 и стиле `\insertbiblioauthorimportant`.
        %  * ни на что не влияют в противном случае
        \nocite{vakbib2}%vak
        \nocite{patbib1}%patent
        \nocite{progbib1}%program
        \nocite{bib1}%other
        \nocite{confbib1}%conf
    \end{refsection}%
        %
        % Всё, что вне этих двух refsection, это refsection=0,
        %  * для диссертации - это нормальные ссылки, попадающие в обычную библиографию
        %  * для автореферата:
        %     * при usefootcite==0, ссылка корректно сработает только для источника из `external.bib`. Для своих работ --- напечатает "[0]" (и даже Warning не вылезет).
        %     * при usefootcite==1, ссылка сработает нормально. В авторской библиографии будут только процитированные в refsection=0 работы.
}
 % Характеристика работы по структуре во введении и в автореферате не отличается (ГОСТ Р 7.0.11, пункты 5.3.1 и 9.2.1), потому её загружаем из одного и того же внешнего файла, предварительно задав форму выделения некоторым параметрам

\underline{\textbf{Объём и структура работы.}} Диссертация состоит из~введения, четырёх глав, заключения и~четырёх приложений. Полный объём диссертации составляет 152~страницы, включая 36~рисунков и~10~таблиц. Список литературы содержит 130~наименований.

\pdfbookmark{Содержание работы}{description}                          % Закладка pdf
\section*{Содержание работы}

Во \underline{\textbf{введении}} обосновывается актуальность темы диссертации, формулируются цель и задачи работы, излагаются научная новизна, практическая значимость, положения, выносимые на защиту, и другие сведения о работе.

\underline{\textbf{Первая глава}} посвящена анализу моделей окружающего пространства БТС и алгоритмов их построения, выбору модели, используемой в работе, и постановке частных задач исследования.

В \underline{\textbf{разделе~1.1}} рассматривается современное состояние беспилотного транспорта и приводится классификация уровней автоматизации SAE~J3016. В \underline{\textbf{разделе~1.2}} описаны применяемые в колёсной робототехнике датчики -- лидары, радары, сонары и камеры -- и ограничения каждого из них: угловое разрешение, дальность, чувствительность к погодным условиям, отсутствие прямого измерения расстояния. Делается вывод о том, что ни один тип датчиков не позволяет построить модель окружающего пространства в одиночку.

В \underline{\textbf{разделе~1.3}} выделены три категории двумерных моделей: объектно-ориентированные представления, сеточные представления и непрерывные представления. Подробно описана карта проходимости пространства, восходящая к работам Х.~Моравека и А.~Элфеса, и байесовский вывод с использованием обратной модели наблюдения как принятый метод обновления её элементов. Каждой клетке $m_i$ сопоставляется логарифм шансов занятости, обновляемый с получением новых данных:

\begin{equation}
	l_{t,i} = l_{t-1,i} + \ln\dfrac{p(m_i|z_t, x_t)}{1 - p(m_i|z_t, x_t)} - l_{0,i},
	\label{eq:syn_log_odd_upd}
\end{equation}
где $l_{0,i}$ -- логарифм шансов, посчитанный для априорной вероятности занятости клетки. \underline{\textbf{Раздел~1.4}} описывает алгоритмы комплексирования данных и три их классификации: по стадии обработки, по используемому математическому аппарату и по характеру взаимодействия комплексируемых данных.

\underline{\textbf{Раздел~1.5}} посвящён метрикам оценки качества моделей. Для объектно-ориентированных представлений рассматривается коэффициент Жаккара. Базовой метрикой сеточных представлений служит среднеквадратичная ошибка относительно эталонной карты

\begin{equation}
	\text{MSE}(M, M_r) = \dfrac{1}{|M|}\sum_{i}\left(p(m_i) - p(m_i^r)\right)^2,
	\label{eq:syn_map_mse}
\end{equation}
где $M$ -- оцениваемая карта, $M_r$ -- эталонная карта того же размера и разрешения, $p(m_i)$ и $p(m_i^r)$ -- вероятности занятости $i$-й клетки, $|M|$ -- число клеток. Эта метрика штрафует ошибку одинаково независимо от того, в какой части карты она допущена. Для планирования пути такое равенство неверно. Карта, на которой препятствие зарегистрировано дальше истинного положения, разрешает движение вперёд и ведёт к столкновению, тогда как карта с той же по величине ошибкой, но сместившей препятствие ближе, приводит лишь к более осторожному объезду. Значения MSE и коэффициента Жаккара у обеих карт совпадают.

Метрика PFC-MSE, предложенная Хорелем, оценивает карту через ту величину, ради которой карта строится, -- через стоимость планируемого по ней пути. Каждой клетке приписывается стоимость проезда $cost(c)$, для занятых клеток существенно большая, чем для свободных (в оригинальной работе $\text{cost}_\text{occ} = 100$ и $\text{cost}_\text{free} = 1$). По карте строится граф, вершинами которого служат клетки, а вес ребра между соседними клетками равен $w(c_1\rightarrow c_2) = cost(c_2)\,\lVert c_2 - c_1\rVert_2$. Алгоритмом Дейкстры от клетки, отвечающей положению БТС, вычисляется стоимость кратчайшего пути $path\_cost(c)$ до каждой клетки, после чего она нормируется на длину пути $L(c)$:

\begin{equation}
	path\_cost_\text{final}(c) = \dfrac{path\_cost(c) - L(c)\,\text{cost}_\text{free}}{\text{cost}_\text{occ} - \text{cost}_\text{free}}.
	\label{eq:syn_path_cost_norm}
\end{equation}

Значения $path\_cost_\text{final}$ по всем клеткам образуют карту стоимостей путей. Такие карты, построенные по эталонной и по оцениваемой картам проходимости, сравниваются среднеквадратичной ошибкой, взвешенной вероятностью того, что клетка свободна:

\begin{equation}
	\text{PFC-MSE}(GT, I) = \dfrac{\sum_c w_c\left(GT_{cost}(c) - I_{cost}(c)\right)^2}{\sum_c w_c},\qquad w_c = p(c = \text{free}),
	\label{eq:syn_pfc_mse}
\end{equation}
где $GT_{cost}$ и $I_{cost}$ -- карты стоимостей путей, построенные по эталонной и по оцениваемой картам проходимости. Ошибка в описании отдельной клетки входит в PFC-MSE через изменение стоимостей путей во все клетки, лежащие за ней. Поэтому ложно свободная клетка, открывающая проезд сквозь препятствие, меняет стоимости в целой области карты, а ложно занятая клетка вплотную к уже известному препятствию не меняет почти ничего. Именно это соответствие веса ошибки её влиянию на планирование и делает метрику пригодной в качестве основной.


В \underline{\textbf{разделе~1.6}} обоснован выбор карты проходимости пространства в качестве базовой модели. Приведены три довода. Модель описывает объекты произвольной геометрии и не требует их отнесения к заранее определённым классам. Она допускает явное представление областей, о проходимости которых сведения отсутствуют. Обновление по формуле~\eqref{eq:syn_log_odd_upd} сводит учёт данных источника к коммутативному сложению вклада в независимую клетку, что снимает необходимость глобальной ассоциации данных между источниками. В качестве метрик выбраны коэффициент Жаккара и PFC-MSE.

\underline{\textbf{Раздел~1.7}} ставит две частные задачи. Первая связана с тем, что камера превосходит остальные бортовые датчики по угловому разрешению и даёт семантическое описание наблюдаемого, однако изображение задаёт лишь направление на объект и не содержит расстояния до него. Учёт визуальных данных поэтому требует проецирования результатов обработки изображения на плоскость движения БТС, а существующие способы такого проецирования тяготеют к двум крайностям: методы, опирающиеся на предположение о плоской поверхности движения, вносят растущую с дальностью погрешность, а обучаемые преобразования признаков изображения в вид сверху требуют производительного графического ускорителя для запуска нейросетевой модели по каждому источнику данных, размеченных мультимодальных наборов для обучения и переобучения при любом изменении состава и расположения датчиков.

Вторая задача связана с ростом задержки $\Delta t$ между регистрацией данных датчиком и их учётом в модели. На неё накладываются два независимых условия. Условие~1 (непотеря данных) требует, чтобы измерения были учтены до поступления следующих: $\Delta t \le 1/f$, что при частоте $f = 10$~Гц даёт $100$~мс. Условие~2 (содержательность разрешения) строже: собственное перемещение БТС за это время скомпенсировано при добавлении фрагмента на карту, однако положение подвижных объектов сцены устаревает. Смещение объекта за время задержки не должно превышать размера элемента карты, иначе разрешение модели оказывается фиктивным, откуда

\begin{equation}
	\Delta t \le \frac{1}{v_{obj}\,res},
	\label{eq:syn_dt_limit}
\end{equation}
где $v_{obj}$ -- собственная скорость объекта, $res$ -- разрешение карты. При $res = 10$~элементов на метр и скорости объекта $40$~км/ч, расчётной для области эксплуатации рассматриваемых БТС -- закрытых территорий и складов открытого типа, -- условие~2 даёт около $9$~мс, тогда как профилирование, приведённое в главе~4, даёт $3,03$~мс при одном источнике данных и $22,78$~мс при восьми: условие~1 выполняется, условие~2 при восьми источниках нарушено. Показано также, что предельное число одновременно учитываемых источников обратно пропорционально четвёртой степени рабочей скорости БТС.

\underline{\textbf{Раздел~1.8}} формулирует требования к разрабатываемым модели окружающего пространства и алгоритмам её построения:

\begin{enumerate}[beginpenalty=10000]
	\item алгоритмы построения модели должны комплексировать данные камер и лидаров, поскольку ни один из этих типов датчиков в отдельности не даёт одновременно высокого углового разрешения и непосредственного измерения расстояния;
	\item алгоритмы не должны опираться на предположение о совпадении поверхности движения с плоскостью;
	\item алгоритмы должны выполняться в режиме реального времени на бортовом вычислителе БТС при ограниченных вычислительных ресурсах, что исключает методы, требующие инференса крупной нейросетевой модели для каждого источника данных;
	\item модель должна быть масштабируемой по числу используемых датчиков и расширяемой под новые их типы;
	\item модель и алгоритм её построения должны допускать одновременное обновление модели несколькими источниками данных.
\end{enumerate}

По результатам первой главы в качестве модели окружающего пространства выбрана карта проходимости, обновляемая байесовским выводом с обратной моделью наблюдения, а в качестве метрик её оценки -- коэффициент Жаккара и PFC-MSE. Поставлены две частные задачи: проецирование визуальных данных без предположения о плоской поверхности движения и одновременное обновление карты несколькими источниками данных. Сформированы требования к используемым моделям и алгоритмам её построения.

\underline{\textbf{Вторая глава}} посвящена алгоритмам проецирования визуальных данных на плоскость движения БТС.

В \underline{\textbf{разделе~2.1}} предложен алгоритм оценки положения окружающих ТС по данным монокулярной камеры. Алгоритм выполняется в два этапа. На первом этапе двумерная детекция уточняется по семантической сегментации проезжей части и проецируется на вид сверху (рисунок~\ref{fig:syn_lshape}). На втором этапе полученная проекция вписывается в ориентированную ограничивающую рамку. Для этого вводится модель L-образной формы видимой границы ТС: вписываемые точки считаются принадлежащими двум смежным ортогональным сторонам объекта, упорядоченными по принадлежности этим сторонам, а разделяющий их индекс в массиве $m \in \{1,\dots,N_p-1\}$ неизвестен заранее. Ранее эта модель применялась лишь к кластерам точек дальномерных датчиков, причём ориентация рамки определялась в ней перебором угла с фиксированным шагом. Поиск рамки сводится к нахождению двух перпендикулярных прямых $l: y = kx + d_1$ и $l': x = -ky + d_2$, минимизирующих сумму квадратов отклонений точек от соответствующей прямой. При фиксированном $m$ задача сводится к линейной регрессии. Показано, что все входящие в решение величины выражаются через семь сумм координат вписываемых точек, каждая из которых пересчитывается при переходе от $m$ к $m+1$ за фиксированное число операций. Это позволяет перебрать всех возможных кандидатов $m$ за $\Theta(N_p)$ операций вместо $\Theta(N_p^2)$ у наивной реализации, где $N_p$ -- число вписываемых точек, и не требует перебора ориентаций рамки.

\begin{figure}[htbp]
	\centering
	\includegraphics[width=0.72\linewidth]{l-shape.png}
	\caption{Проекция двумерной ограничивающей рамки на вид сверху с учётом
	сегментации проезжей части.}
	\label{fig:syn_lshape}
\end{figure}

\underline{\textbf{Раздел~2.2}} описывает исследование точности предложенного алгоритма вписывания. Он сравнивался с построением рамки минимальной площади и четырьмя алгоритмами вписывания, определяющими ориентацию рамки по выпуклой оболочке вписываемых точек: Rotating PCA, Diagonal PCA, Longest Diagonal и Rotated Triangle. Все четыре изначально разработаны для кластеров точек дальномерных датчиков и здесь применены к тем же входным данным, что и предложенный алгоритм. На записях показаний датчиков реальных БТС предложенный алгоритм показал наивысший средний коэффициент Жаккара $0,337$ против $0,328$ у лучшего из сравниваемых.

\begin{table}[htbp]
	\small\centering
	\caption{Коэффициент Жаккара для алгоритмов вписывания проекции объекта на виде сверху в ориентированную ограничивающую рамку. Наилучшие значения выделены полужирным.}
	\label{tab:syn_lshape_res}
	\begin{tabular}{|l|c|c|c|c|}
		\hline
		\textbf{Алгоритм} & \textbf{Среднее} & \textbf{25\,\%} & \textbf{50\,\%} & \textbf{75\,\%} \\
		\hline
		Рамка минимальной площади & $0,294$ & $0,198$ & $0,257$ & $0,383$ \\
		Rotating PCA & $0,328$ & $0,219$ & $0,316$ & $0,448$ \\
		Diagonal PCA & $0,321$ & $0,202$ & $0,316$ & $0,445$ \\
		Longest Diagonal & $0,323$ & $0,199$ & $0,287$ & $0,442$ \\
		Rotated Triangle & $0,301$ & $0,208$ & $0,264$ & $0,425$ \\
		\textbf{L-Shape} & \boldmath$0,337$ & \boldmath$0,230$ & \boldmath$0,328$ & \boldmath$0,451$ \\
		\hline
	\end{tabular}
\end{table}


В \underline{\textbf{разделе~2.3}} предложен алгоритм оценки положения окружающих ТС по комплексированию визуальных и лидарных данных. Ключевым его шагом является разделение облака точек на точки земли и точки препятствий, выполняемое без предположения о плоской либо кусочно-плоской поверхности движения. Вместо поиска плоскостей высота поверхности $z_l$ оценивается регрессией на гауссовских процессах по координатам $(x_l, y_l)$. В качестве ядра выбрано ядро Матерна при $\nu = 3/2$, лучше приближающее функцию в точках разрыва. Точка считается не принадлежащей земле, если её измеренная координата превышает предсказанную на порог, зависящий от расстояния до БТС. Результат работы алгоритма приведён на рисунке~\ref{fig:syn_splitting}.

\begin{figure}[htbp]
	\centering
	\includegraphics[width=0.6\linewidth]{splitting_demo_pc.png}
	\caption{Пример работы алгоритма разделения облака точек на точки земли
	и препятствий.}
	\label{fig:syn_splitting}
\end{figure}

\underline{\textbf{Раздел~2.4}} содержит исследование точности алгоритма комплексирования. Сравнение проводилось между сквозными конвейерами -- от сырых данных кадра до ориентированной ограничивающей рамки -- на общем наборе эталонных данных, по единой метрике и при едином сопоставлении вычисленных и эталонных рамок венгерским алгоритмом. Предложенный алгоритм показал средний коэффициент Жаккара $0,439$ против $0,399$ у лучшего из алгоритмов, использующих только лидарные данные, и $0,337$ у алгоритма, использующего только визуальные данные, что соответствует приросту на $10,0$~\% и $30,3$~\% соответственно. Среднее время обработки одного кадра составило $66,7$~мс, что удовлетворяет условию~1 на задержку ($100$~мс при частоте регистрации данных $10$~Гц).

\begin{table}[htbp]
	\small\centering
	\caption{Средний коэффициент Жаккара для сквозных конвейеров оценки положения окружающих ТС. Сравниваются конвейеры целиком, от сырых данных кадра до ориентированной рамки.}
	\label{tab:syn_fusion_res}
	\begin{tabular}{|l|c|c|}
		\hline
		\textbf{Конвейер} & \textbf{Данные} & \textbf{Среднее} \\
		\hline
		Rotating PCA & лидар & $0,396$ \\
		Diagonal PCA & лидар & $0,399$ \\
		Longest Diagonal & лидар & $0,335$ \\
		Rotated Triangle & лидар & $0,339$ \\
		L-Shape & камера & $0,337$ \\
		\textbf{Предложенный} & лидар и камера & \boldmath$0,439$ \\
		\hline
	\end{tabular}
\end{table}


В \underline{\textbf{разделе~2.5}} предложен алгоритм проекции семантической сегментации проезжей части на плоскость движения БТС. Проецирование выполняется не на фиксированную плоскость $z_l = 0$, а на поверхность $z = z(x, y)$, значения которой в центрах элементов сетки оцениваются той же регрессией на гауссовских процессах (рисунок~\ref{fig:syn_proj}). Полученная проекция объединяется с картой базового лидарного алгоритма конкурирующим комплексированием: элементы, о проходимости которых источники дают противоречивые сведения, переводятся в состояние <<неизвестно>>.

\begin{figure}[htbp]
	\centering
	\includegraphics[width=0.49\linewidth]{PlaneProj.png}
	\hfill
	\includegraphics[width=0.49\linewidth]{NonPlaneProj.png}
	\caption{Проецирование сегментации проезжей части на плоскость $z_l = 0$ (слева)
	и на оценённую поверхность $z = z(x,y)$ (справа).}
	\label{fig:syn_proj}
\end{figure}

\underline{\textbf{Раздел~2.6}} описывает исследование точности полученного объединения алгоритмов. Использовались пять записей показаний датчиков БТС, каждая из которых содержит объединённое облако точек с двух лидаров и изображения с трёх бортовых камер. Эталонные карты имели размер $600\times600$ элементов при разрешении $10$~элементов на метр, дальность проецирования ограничивалась $40$~м. По сравнению с картой, строящейся исключительно по лидарным данным, среднеквадратичная ошибка относительно эталонной карты снизилась в среднем с $0,231$ до $0,218$, а метрика PFC-MSE -- с $1019,3$ до $924,4$, причём по основной метрике улучшение достигнуто во всех пяти сценариях. Примеры построенных карт приведены на рисунке~\ref{fig:syn_maps}.

Абляционное исследование раздельно оценивает вклад двух составляющих алгоритма. Вариант без учёта рельефа, проецирующий сегментацию на фиксированную плоскость, уступает по всем трём метрикам и оказывается хуже базового лидарного алгоритма, вовсе не использующего визуальные данные. Вариант с дополняющим комплексированием вместо конкурирующего даёт наилучшую среднеквадратичную ошибку $0,183$, одновременно оказываясь худшим по PFC-MSE $1130,8$: близость карты к эталонной по поэлементному сходству не влечёт её пригодности для планирования маршрута. Среднее время одной итерации построения карты составило $1,63$~мс.

\begin{table}[htbp]
	\small\centering
	\caption{Усреднённые по пяти сценариям значения метрик: Б~-- базовый лидарный алгоритм, П~-- предложенный, А1~-- проецирование на фиксированную плоскость, А2~-- дополняющее комплексирование вместо конкурирующего. Наилучшие значения выделены полужирным.}
	\label{tab:syn_ablation_res}
	\begin{tabular}{|l|c|c|c|}
		\hline
		\textbf{Алгоритм} & \textbf{СКО} \tiny$\downarrow$ & \boldmath$mJ$ \tiny$\uparrow$ & \textbf{PFC-MSE} \tiny$\downarrow$ \\
		\hline
		Б & $0,231$ & $0,225$ & $1019,3$ \\
		\textbf{П} & $0,218$ & \boldmath$0,238$ & \boldmath$924,4$ \\
		А1 & $0,236$ & $0,230$ & $1059,4$ \\
		А2 & \boldmath$0,183$ & $0,205$ & $1130,8$ \\
		\hline
	\end{tabular}
\end{table}


\begin{figure}[htbp]
	\centering
	\includegraphics[width=0.72\linewidth]{table_with_maps_new.png}
	\caption{Эталонные карты проходимости пространства и карты, построенные
	сравниваемыми алгоритмами на тестовых данных.}
	\label{fig:syn_maps}
\end{figure}

Итогом второй главы стали три алгоритма проецирования визуальных данных на плоскость движения БТС. На записях показаний датчиков реальных БТС показано, что отказ от приближения поверхности движения плоскостью повышает и точность оценки положения окружающих ТС (коэффициент Жаккара $0,439$ против $0,399$ у лучшего лидарного алгоритма), и качество карты проходимости пространства (PFC-MSE $924,4$ против $1019,3$).

\underline{\textbf{Третья глава}} посвящена программному комплексу построения карты проходимости пространства Occupancy Map Fusion, реализующему предложенные алгоритмы.

Комплекс должен удовлетворять требованиям масштабируемости по числу датчиков, расширяемости под новые их типы и возможности одновременного обновления модели несколькими источниками, сформулированным в первой главе. \underline{\textbf{Раздел~3.3}} описывает принятые для этого архитектурные решения. Каждый источник данных представлен поставщиком фрагмента карты: преобразование его данных во фрагмент выполняется через единый интерфейс, и подключение датчика нового типа сводится к реализации одного метода, не требуя изменения кода построения карты. Класс, строящий карту, не знает о конкретных поставщиках и получает фрагменты через общую очередь, поэтому число одновременно подключённых источников не фиксировано в коде. Значения элементов фрагмента вычисляются по задаваемым для каждого поставщика оценкам вероятностей верного срабатывания и пропуска препятствия. Диаграмма компонентов комплекса приведена на рисунке~\ref{fig:syn_omf}.

Комплекс реализован на языках C++ и CUDA~C++ и обеспечивает многопоточное обновление карты по данным пяти типов источников, включая комплексирование двумерных детекций и семантической сегментации проезжей части с лидарным облаком. Предусмотрено подключение и отключение источников по ходу выполнения программы, а также опциональное удаление из карты областей, занимаемых динамическими препятствиями.

\begin{figure}[htbp]
	\centering
	\includegraphics[width=0.7\linewidth]{OMF_UML.png}
	\caption{Диаграмма компонентов программного комплекса Occupancy Map Fusion
	на языке UML.}
	\label{fig:syn_omf}
\end{figure}

Третья глава даёт программный комплекс, реализующий предложенные алгоритмы и удовлетворяющий сформулированным в первой главе требованиям: подключение источника нового типа сводится к реализации одного преобразования данных во фрагмент карты, а число одновременно подключённых источников не ограничено архитектурой комплекса.

\underline{\textbf{Четвёртая глава}} посвящена параллельному алгоритму построения карты проходимости пространства.

В \underline{\textbf{разделе~4.1}} приведены результаты профилирования комплекса при одном и восьми одновременно подключённых источниках данных (рисунок~\ref{fig:syn_perf}). Время обновления карты составило $3,03$ и $22,78$~мс соответственно. В обоих случаях удовлетворяется условие~1 ($100$~мс при частоте регистрации данных $10$~Гц), однако при восьми источниках допустимая по условию~2 задержка $9$~мс превышена примерно в $2,5$ раза. Рост задержки более чем в семь раз показывает, что последовательное добавление фрагментов ограничивает дальнейшее наращивание числа источников.

\begin{figure}[htbp]
	\centering
	\includegraphics[width=0.8\linewidth]{perf-omf-8.png}
	\caption{Задержки обновления карты проходимости пространства при восьми
	источниках данных.}
	\label{fig:syn_perf}
\end{figure}

\underline{\textbf{Раздел~4.2}} формализует задачу построения локальной карты. Каждый источник независимо от остальных порождает фрагменты карты -- двумерные массивы логарифмов шансов занятости с известными положением и временем регистрации, -- которые требуется свести в общую карту, следующую за БТС. Над картой выполняются три операции: сдвиг вслед за перемещением БТС, добавление фрагмента и получение карты для передачи потребителям. В базовой реализации каждая из них требует прохода по всем элементам карты и выполняется под исключающей блокировкой, что и порождает наблюдаемую сериализацию.

\underline{\textbf{Раздел~4.3}} предлагает три последовательные модификации базовой реализации. Каждая следующая снимает то ограничение на одновременный доступ, которое осталось от предыдущей.

\begin{enumerate}[beginpenalty=10000]
	\item \emph{Расширенная карта} направлена на стоимость сдвига. Массив хранит область больше отображаемой, а положение видимой части задаётся отдельной координатой. Сдвиг сводится к изменению этой координаты и перемещает содержимое лишь изредка, когда видимая область подходит к границе массива, поэтому его амортизированная сложность составляет $\Theta(1)$ вместо $\Theta(N)$. Все операции при этом остаются блокирующими.
	\item \emph{Карта с разделяемыми блокировками} изменяет структуру самой клетки: вместо одного числа клетка хранит атомарную пару <<логарифм шансов -- время последнего обновления>>. Требование, чтобы вся карта относилась к одному моменту времени, тем самым снимается, и приведение всех клеток к общему времени становится ненужным. Клетка обновляется операцией сравнения с обменом, поэтому источники, фрагменты которых не пересекаются, добавляют их одновременно. Исключительные права на карту нужны теперь только сдвигу.
	\item \emph{Карта с полным обновлением без блокировок} снимает и это ограничение. Карта разбивается на два примыкающих фрагмента, хранимых в атомарных разделяемых указателях. Сдвиг заменяет тот фрагмент, на котором видимая область в данный момент не находится, и потому выполняется, не блокируя ни добавление фрагментов, ни получение карты.
\end{enumerate}

\underline{\textbf{Раздел~4.4}} содержит сравнение задержки обновления карты. Замеры выполнялись средствами библиотеки Google Benchmark для видимых областей размером от $10$ до $4000$ элементов. Каждый замер повторялся десять раз, коэффициент вариации по повторам не превышает $2,5$~\%. Наиболее показательным является сценарий цикла публикации карты, в котором в восьми потоках выполняются сдвиг карты и добавление фрагмента, после чего карта извлекается для передачи потребителям. Все реализации имеют линейную вычислительную сложность относительно размера карты, и каждая следующая модификация превосходит предыдущую. Лучший результат показала карта с полным обновлением без блокировок: коэффициент линейной аппроксимации зависимости времени итерации от размера карты у неё в $13,5$ раза меньше, чем у базовой реализации. Для карты с разделяемыми блокировками отношение составляет $10,3$, для расширенной карты -- $6,2$ (рисунок~\ref{fig:syn_bench}).

\begin{table}[htbp]
	\small\centering
	\caption{Во сколько раз коэффициент линейной аппроксимации времени итерации цикла публикации карты меньше, чем у исходной реализации.}
	\label{tab:syn_bench_res}
	\begin{tabular}{|l|c|}
		\hline
		\textbf{Реализация} & \textbf{Отношение к исходной} \\
		\hline
		Расширенная карта & $6,2$ \\
		С разделяемыми блокировками & $10,3$ \\
		\textbf{С полным обновлением без блокировок} & \boldmath$13,5$ \\
		\hline
	\end{tabular}
\end{table}


\begin{figure}[htbp]
	\centering
	\includegraphics[width=0.75\linewidth]{_SinglePub_it4.png}
	\caption{Зависимость времени одной итерации цикла публикации карты от размера
	карты для базовой реализации и трёх предложенных модификаций.}
	\label{fig:syn_bench}
\end{figure}

В \underline{\textbf{разделе~4.5}} представлено расширение карты с полным обновлением без блокировок на двумерный случай. Полная карта окружения разбивается на девять двумерных массивов атомарных пар, и при выходе видимой области за отмеченные границы создаются новые массивы, замещающие часть прежних так, чтобы массив с видимой областью стал центральным (рисунок~\ref{fig:syn_2d}). Определение условий, при которых такая реализация эффективнее исходной, составляет предмет дальнейшей работы.

\begin{figure}[htbp]
	\centering
	\includegraphics[width=0.55\linewidth]{Lockfree2D.png}
	\caption{Двумерная реализация локальной карты проходимости пространства
	с полным обновлением без блокировок.}
	\label{fig:syn_2d}
\end{figure}


В четвёртой главе получены структура данных и алгоритм обновления карты, в которых операции над разделяемой сеткой выполняются без примитивов взаимного исключения. Коэффициент линейной аппроксимации времени итерации цикла публикации карты уменьшен в $13,5$ раза относительно исходной реализации.

\FloatBarrier
\pdfbookmark{Заключение}{conclusion}                                  % Закладка pdf
В \underline{\textbf{заключении}} приведены основные результаты работы, которые заключаются в следующем:
%% Согласно ГОСТ Р 7.0.11-2011:
%% 5.3.3 В заключении диссертации излагают итоги выполненного исследования, рекомендации, перспективы дальнейшей разработки темы.
%% 9.2.3 В заключении автореферата диссертации излагают итоги данного исследования, рекомендации и перспективы дальнейшей разработки темы.
%% Поэтому эта часть общая и загружается из одного файла в автореферат и в диссертацию.
%%
%% ВНИМАНИЕ: файл подключается и диссертацией (Dissertation/conclusion.tex), и авторефератом
%% (Synopsis/content.tex). Ссылки \cref/\ref здесь использовать нельзя: метки глав и разделов
%% существуют только в диссертации, в автореферате они дадут неразрешённые ссылки.
%% Указания на главы и разделы приводятся текстом.
\begin{enumerate}
	\item Выполнен анализ моделей окружающего пространства БТС, алгоритмов их построения и способов комплексирования разнородных данных. На его основе обоснован выбор карты проходимости пространства в качестве модели, описывающей проходимость окружающего БТС пространства, байесовского вывода с использованием обратной модели наблюдения в качестве алгоритма её построения, а также коэффициента Жаккара и метрики PFC-MSE в качестве метрик оценки её качества. Сформулированы требования к разрабатываемым модели и алгоритмам её построения.
	\item Предложен алгоритм оценки положения окружающих ТС на основе комплексирования визуальных и лидарных данных, выделяющий точки поверхности движения оценкой её формы регрессией на гауссовских процессах. На реальных данных алгоритм показал средний коэффициент Жаккара $0,439$ против $0,399$ у лучшего лидарного и $0,337$ у визуального алгоритмов, не использующих обучаемых преобразований в вид сверху. Алгоритм сохраняет превосходство во всём диапазоне рассмотренных дальностей.
	\item Предложен алгоритм проецирования семантической сегментации проезжей части с изображения на оценённую по лидарным данным поверхность движения, а также способ его объединения с базовым алгоритмом построения карты проходимости пространства, использующий конкурирующее комплексирование. Результаты экспериментов показали, что полученное объединение алгоритмов снижает среднеквадратичную ошибку стоимости поиска пути (PathFinding Cost Mean Squared Error, PFC-MSE) с $1019,3$ до $924,4$ в пяти рассмотренных сценариях, а среднеквадратичную ошибку относительно эталонной карты -- с $0,231$ до $0,218$. Проецирование той же сегментации на фиксированную плоскость увеличивает PFC-MSE до $1059,4$. 
	\item Разработан программный комплекс построения карты проходимости пространства на основе комплексирования разнородных данных от датчиков, установленных на БТС, и их обработчиков. Архитектура программного комплекса позволяет осуществлять параллельную подготовку фрагментов карты по показаниям независимых источников данных, оставляя последовательным только этап добавления фрагментов в общую карту.
	\item Реализованы три последовательные модификации структуры данных для представления карты проходимости пространства, направленные на снижение вычислительной сложности производимых операций. Тестирование на синтетических данных, моделирующих одномерную карту проходимости пространства, показало, что наибольший выигрыш для одного цикла публикации карты даёт финальная модификация -- карта проходимости пространства с полным обновлением без блокировок: коэффициент линейной аппроксимации зависимости времени итерации от размера карты уменьшается в $13,5$ раза относительно исходной реализации.
	\item Предложено расширение последней модификации карты проходимости пространства для двумерного пространства.
\end{enumerate}

По результатам работы могут быть даны следующие рекомендации. При построении карты проходимости пространства на БТС, оснащённых камерами и лидарами, проецирование результатов обработки изображения целесообразно выполнять на оценённую по дальномерным данным поверхность движения, а не на фиксированную плоскость: как показано во второй главе, отказ от предположения о плоской поверхности улучшает описание сцены даже по сравнению с алгоритмом, вовсе не использующим визуальные данные. Объединение фрагментов карты, построенных по разным источникам данных, следует выполнять конкурирующим комплексированием, переводя элементы карты, о проходимости которых источники дают противоречивые сведения, в состояние <<неизвестно>>. Наконец, при проектировании программного обеспечения построения карты источники данных целесообразно подключать через единый интерфейс поставщика фрагментов, что позволяет наращивать их число без изменения кода построения карты.

Перспективы дальнейшей разработки темы состоят в следующем. Двумерная реализация карты проходимости пространства с полным обновлением без блокировок требует переноса и валидации на реальной платформе, а также измерения задержки обновления карты при различном числе источников данных в тех же условиях, в которых выполнено профилирование программного комплекса. Это позволит определить условия, при которых она вычислительно эффективнее исходной реализации. Разработанный программный комплекс допускает подключение поставщиков фрагментов карты по данным радаров и сонаров, однако качество получаемой при этом модели в настоящей работе не исследовалось. Отдельным направлением является переход к динамической карте проходимости пространства, хранящей вместе с занятостью элемента оценку скорости наблюдаемого препятствия, что позволило бы отказаться от механизма затухания как основного способа учёта подвижных объектов. Отдельного внимания заслуживает используемая в работе метрика PFC-MSE: стоимости путей в ней вычисляются без учёта кинематической модели движения и габаритов БТС, из-за чего области, недостижимые для БТС, дают вклад в итоговое значение метрики наравне с достижимыми. Учёт кинематической модели движения и габаритов БТС при построении графа сделал бы оценку влияния ошибок карты на алгоритмы планирования пути ближе к тому, как карта используется в действительности. Наконец, представляет интерес оценка предложенных алгоритмов на публичных наборах данных: их применение потребует либо подбора наборов с сопоставимой конфигурацией датчиков, либо построения эталонной разметки проходимости пространства по имеющимся в них аннотациям.


\pdfbookmark{Литература}{bibliography}                                % Закладка pdf

\ifdefmacro{\microtypesetup}{\microtypesetup{protrusion=false}}{} % не рекомендуется применять пакет микротипографики к автоматически генерируемому списку литературы
\urlstyle{rm}                               % ссылки URL обычным шрифтом
%
% Работы автора цитируются в порядке групп: ВАК, издания, индексируемые Web of Science
% и Scopus, патенты, программы для ЭВМ. Порядок \nocite определяет порядок в библиографии.
\nocite{Chekanov_sensys_2025}%
\nocite{Voronkov_ip_2026}%
\nocite{чеканов2026многопоточные}%
\nocite{Chekanov_icmv_2024}%
\nocite{Chekanov_ecms_2026}%
\nocite{patbib1}%
\nocite{progbib1}%
\nocite{progbib2}%
%
% Группировка списка работ автора выполняется вручную, а не командой
% \insertbiblioauthorgrouped: две работы помечены одновременно authorwos и authorscopus,
% и стандартная команда, печатающая эти группы раздельно, вывела бы их дважды.
% Объединённая группа задаётся фильтром scopuswos, уже определённым в biblio/biblatex.tex.
\begin{refcontext}[labelprefix={}]
    \section*{\bibtitleauthor}
    \printbibliography[heading=pubsubgroup, section=0, keyword=biblioauthorvak,
        title=\bibtitleauthorvak, resetnumbers=true]
    \printbibliography[heading=pubsubgroup, section=0, filter=scopuswos,
        title={В изданиях, индексируемых Web of Science и Scopus}, resetnumbers=false]
    \printbibliography[heading=pubsubgroup, section=0, keyword=biblioauthorpatent,
        title=\bibtitleauthorpatent, resetnumbers=false]
    \printbibliography[heading=pubsubgroup, section=0, keyword=biblioauthorprogram,
        title=\bibtitleauthorprogram, resetnumbers=false]
\end{refcontext}
\ifdefmacro{\microtypesetup}{\microtypesetup{protrusion=true}}{}
\urlstyle{tt}                               % возвращаем установки шрифта ссылок URL
